%; whizzy chapter
% -initex iniptex -latex platex -format platex -bibtex jbibtex -fmt fmt
% 以上 whizzytex を使用する場合の設定。

%     Kansai Debian Meeting resources
%     Copyright (C) 2007 Takaya Yamashita
%     Thank you for Tokyo Debian Meeting resources

%     This program is free software; you can redistribute it and/or modify
%     it under the terms of the GNU General Public License as published by
%     the Free Software Foundation; either version 2 of the License, or
%     (at your option) any later version.

%     This program is distributed in the hope that it will be useful,
%     but WITHOUT ANY WARRANTY; without even the implied warranty of
%     MERCHANTABILITY or FITNESS FOR A PARTICULAR PURPOSE.  See the
%     GNU General Public License for more details.

%     You should have received a copy of the GNU General Public License
%     along with this program; if not, write to the Free Software
%     Foundation, Inc., 51 Franklin St, Fifth Floor, Boston, MA  02110-1301 USA

%  preview (shell-command (concat "evince " (replace-regexp-in-string "tex$" "pdf"(buffer-file-name)) "&"))
% 画像ファイルを処理するためにはebbを利用してboundingboxを作成。
%(shell-command "cd image200708; ebb *.png")

%%ここからヘッダ開始。

\documentclass[mingoth,a4paper]{jsarticle}
\usepackage{kansaimonthlyreport}
\usepackage[dvipdfmx]{xy}
\usepackage{etex}
\usepackage{ulem}

% 日付を定義する、毎月変わります。
\newcommand{\debmtgyear}{2018}
\newcommand{\debmtgdate}{25}
\newcommand{\debmtgmonth}{2}
\newcommand{\debmtgnumber}{132}

\def\fixme#1{{\color{red}{#1}}}

\begin{document}

\begin{titlepage}

% 毎月変更する部分、本文の末尾も修正することをわすれずに

  第\debmtgnumber{}回 関西 Debian 勉強会資料

  \vspace{2cm}

  \begin{center}
    \includegraphics{image200802/kansaidebianlogo.png}
  \end{center}

  \begin{flushright}
    \hfill{}関西Debian勉強会担当者 かわだ \\
    \hfill{}\debmtgyear{}年\debmtgmonth{}月\debmtgdate{}日
  \end{flushright}

  \thispagestyle{empty}
\end{titlepage}

\dancersection{Introduction}{Debian JP}

\vspace{1em}

関西Debian勉強会はDebian GNU/Linuxのさまざまなトピック
(新しいパッケージ、Debian特有の機能の仕組、Debian界隈で起こった出来事、
などなど）について話し合う会です。

目的として次の三つを考えています。
\begin{itemize}
\item MLや掲示板ではなく、直接顔を合わせる事での情報交換の促進
\item 定期的に集まれる場所
\item 資料の作成
\end{itemize}

それでは、楽しい一時をお過ごしください。

\newpage

\begin{minipage}[b]{0.2\hsize}
  {\rotatebox{90}{\fontsize{80}{80}
{\gt{関西 Debian 勉強会}}}}
\end{minipage}
\begin{minipage}[b]{0.8\hsize}
\hrule
\vspace{2mm}
\hrule
\setcounter{tocdepth}{1}
\tableofcontents
\vspace{2mm}
\hrule
\end{minipage}

\dancersection{最近のDebian関係のイベント報告}{Debian JP}

\subsection{第131回関西Debian勉強会}

131回目の関西Debian勉強会は、1月28日(日)に港区民センターで日本openSUSEユーザ会と
LILO\&東海道らぐと合同でライトニングトーク大会として開催されました。

\subsection{第159回東京エリアDebian勉強会}

159回目の東京エリアDebian勉強会は1月20日(土)に株式会社朝日ネットさんで開催されま
した。

\subsection{第160回東京エリアDebian勉強会}

160回目の東京エリアDebian勉強会は2月24日(土)に明星大学で開催された OSC 2018
Tokyo/Spring に出張されました。

セッションは杉本さんによる「apache に django の web アプリでデプロイしてみよう」
という django を題材として Debian の apache 設定の解説でした。


\subsection{Debian Project}

\subsubsection{Spectre / Meltdown}

Spectre (CVE-2017-5715, CVE-2017-5753) と Meltdown (CVE-2017-5754) と呼ばれてい
る脆弱性への Linux カーネルセキュリティアップデート
\footnote{\url{https://www.debian.org/security/2018/dsa-4120}}
が出ました。

Spectre variant 2 には ``retpoline'' を使った対策が取られています。この対策のた
め gcc-6 に ``retpoline'' 対応が入り、こちらもセキュリティアップデートという形で
更新\footnote{\url{https://www.debian.org/security/2018/dsa-4121}}が出ています。

どちらのパッケージも忘れずに更新しましょう。


\subsubsection{Upcoming stable point release (9.4)}

strech のポイントリリース 9.4 が 3 月 10 日(土)に予定されています
\footnote{\url{https://lists.debian.org/debian-project/2018/02/msg00014.html}}。

6 月で stretch リリースから一年となり jessie のセキュリティサポートも終わります。
みなさん stretch への移行は順調に進んでいますでしょうか。


\subsubsection{Extended Long Term Support for Wheezy}

Debian 7 Wheezy の Long Term Support は 2018 年 5 月を持って終了する予定です。
しかし、ここにきてさらに 10 年の長期サポートを行なう動きがあります
\footnote{\url{https://lists.debian.org/debian-devel/2018/02/msg00470.html}}。

インフラやトラッカー、サインキーなどをどうするかといった課題がありどうなるかはわ
かりませんが、リリースから 15 年となると、サポート期間が最長のディストリビューショ
ンになりそうです。

\subsubsection{CUPS GPL → Apache license change, how to proceed?}

Apple が CUPS のラインセンスを GPL-2.0 with AOSDL exception から Apache Licence
2.0 に変更しました。ライセンス変更された CUPS 2.3 が experimental に入りましたが、
ライセンス不整合を起こすため RC バグが取れるまで先は長そうです
\footnote{\url{https://lists.debian.org/debian-devel/2018/02/msg00227.html}}。


\dancersection{事前課題}{関西Debian勉強会}

今回の事前課題はありませんでした。

参加者の皆さんは以下の通りです:

\begin{prework}{ znz }
\end{prework}

\begin{prework}{ gdevmjc }
\end{prework}

\begin{prework}{ t3rkwd }
\end{prework}

\begin{prework}{ sato\_makoto }
\end{prework}

\begin{prework}{ Katsuki Kobayashi }
\end{prework}

\begin{prework}{ takeuchi }
\end{prework}

\begin{prework}{ satomicchy }
\end{prework}

\begin{prework}{ uwabami }
\end{prework}

\begin{prework}{ ipv6waterstar }
\end{prework}

% ---------------------------------------------------
\dancersection{ufw 再入門}{西山和広}

ufw はバージョンアップに伴い、機能が追加されてできることが増えています。
そこで、現在の Debian の stable である stretch に入っている ufw 0.35-4 を元に基本的な機能やちょっと複雑なネットワークでの使用方法を紹介します。

\subsection{ufw とは?}

iptables のラッパーのようなもので、ファイアウォールの設定を簡単にできるようにするためのものです。
ちなみに名前の ufw は Uncomplicated Firewall の略で、Ubuntu Firewall ではありません。

\subsection{使用開始}

まず「ufw enable」で有効にします。 (「ufw disable」で戻せます。)
有効にするとすぐに反映されて、次回起動時からも有効になります。
ssh で接続している場合は、「ufw enable」の前に、そのポートを許可しておくと良いでしょう。
最近は大丈夫のようですが、昔は enable した直後に反応がなくなって切れてしまうことがあったので、再接続できるように許可しておいたり、他の接続手段も確保しておいたりしておく方が安全です。

\begin{verbatim}
$ sudo ufw allow 22/tcp
Rules updated
Rules updated (v6)
$ sudo ufw enable
Command may disrupt existing ssh connections. Proceed with operation (y|n)? y
Firewall is active and enabled on system startup
\end{verbatim}

\subsection{デフォルトの挙動}

デフォルトでは外向きは許可、内向きは拒否という一般的に推奨される構成になっています。
許可などの設定は IP アドレスなどの指定がない場合、IPv4 と IPv6 の両方を同時に設定できるので、iptables と ip6tables を別々に設定して二度手間になるようなことがありません。

\subsection{基本的な使い方}

ポート番号のみを指定すると TCP/UDP 両方許可できます。
例えば DNS なら、以下のようになります。

\begin{verbatim}
$ sudo ufw allow 53
Rule added
Rule added (v6)
\end{verbatim}

TCP のみ許可するには「ポート番号/tcp」と指定します。
例えば http なら、以下のようになります。

\begin{verbatim}
$ sudo ufw allow 80/tcp
Rule added
Rule added (v6)
\end{verbatim}

UDP のみ許可するには「ポート番号/udp」と指定します。
ポート番号には「:」区切りで範囲を指定できます。
例えば mosh で使われる範囲を許可するなら、以下のようになります。

\begin{verbatim}
$ sudo ufw allow 60000:61000/udp
Rule added
Rule added (v6)
\end{verbatim}

連続しないポート番号を「,」区切りでまとめて指定することもできます。
(「,」で区切られた中に範囲を指定することもできます。)
例えば SMTP などを許可するなら、以下のようになります。

\begin{verbatim}
$ sudo ufw allow 25,465,587/tcp
Rule added
Rule added (v6)
\end{verbatim}

\subsection{設定確認}

「ufw status」で設定を確認できます。
「ufw status verbose」でデフォルトポリシーなどを含めた設定の確認ができます。
「ufw status numbered」で delete や insert で使う番号が確認できます。

\begin{verbatim}
$ sudo ufw status
Status: active

To                         Action      From
--                         ------      ----
22/tcp                     ALLOW       Anywhere
53                         ALLOW       Anywhere
80/tcp                     ALLOW       Anywhere
60000:61000/udp            ALLOW       Anywhere
25,465,587/tcp             ALLOW       Anywhere
22/tcp (v6)                ALLOW       Anywhere (v6)
53 (v6)                    ALLOW       Anywhere (v6)
80/tcp (v6)                ALLOW       Anywhere (v6)
60000:61000/udp (v6)       ALLOW       Anywhere (v6)
25,465,587/tcp (v6)        ALLOW       Anywhere (v6)

$ sudo ufw status verbose
Status: active
Logging: on (low)
Default: deny (incoming), allow (outgoing), disabled (routed)
New profiles: skip

To                         Action      From
--                         ------      ----
22/tcp                     ALLOW IN    Anywhere
53                         ALLOW IN    Anywhere
80/tcp                     ALLOW IN    Anywhere
60000:61000/udp            ALLOW IN    Anywhere
25,465,587/tcp             ALLOW IN    Anywhere
22/tcp (v6)                ALLOW IN    Anywhere (v6)
53 (v6)                    ALLOW IN    Anywhere (v6)
80/tcp (v6)                ALLOW IN    Anywhere (v6)
60000:61000/udp (v6)       ALLOW IN    Anywhere (v6)
25,465,587/tcp (v6)        ALLOW IN    Anywhere (v6)

$ sudo ufw status numbered
Status: active

     To                         Action      From
     --                         ------      ----
[ 1] 22/tcp                     ALLOW IN    Anywhere
[ 2] 53                         ALLOW IN    Anywhere
[ 3] 80/tcp                     ALLOW IN    Anywhere
[ 4] 60000:61000/udp            ALLOW IN    Anywhere
[ 5] 25,465,587/tcp             ALLOW IN    Anywhere
[ 6] 22/tcp (v6)                ALLOW IN    Anywhere (v6)
[ 7] 53 (v6)                    ALLOW IN    Anywhere (v6)
[ 8] 80/tcp (v6)                ALLOW IN    Anywhere (v6)
[ 9] 60000:61000/udp (v6)       ALLOW IN    Anywhere (v6)
[10] 25,465,587/tcp (v6)        ALLOW IN    Anywhere (v6)
\end{verbatim}

\subsection{ルールの削除}

「ufw delete 追加した時のルール」や「ufw delete ルール番号」でルールを削除できます。

\begin{verbatim}
$ sudo ufw delete allow 25,465,587/tcp
Rule deleted
Rule deleted (v6)
$ sudo ufw delete allow 60000:61000/udp
Rule deleted
Rule deleted (v6)
\end{verbatim}

\subsection{ルールの表記方法}

今まで使ってきた

\begin{verbatim}
ufw  [--dry-run] [delete] [insert NUM] allow|deny|reject|limit
     [in|out] [log|log-all]
     [ PORT[/PROTOCOL] | APP‐NAME ]
     [comment COMMENT]
\end{verbatim}

という省略記法の他に、

\begin{verbatim}
ufw [--dry-run] [rule] [delete] [insert NUM] allow|deny|reject|limit
    [in|out  [on  INTERFACE]]  [log|log-all] [proto  PROTOCOL]
    [from  ADDRESS [port PORT | app APPNAME ]]
    [to ADDRESS [port PORT | app APPNAME ]]
    [comment COMMENT]
\end{verbatim}

という細かく指定する書き方もあります。
IP アドレスを指定するには後者の表記方法を使う必要があります。
詳細は man を参照してください。

\subsection{特定 IP アドレスのみ拒否}

前述のように IP アドレスを指定するには「ポート番号/プロトコル」という指定の仕方では無理なので、細かく指定していきます。
また、許可より前に拒否するルールがないといけないので、「insert」 を使います。ここで先ほどの「ufw status numbered」の出力をみると 3 より前にあれば良いとわかるので、「insert 3」と指定します。
ポート番号を指定するときは IP アドレスの指定は必須なので、特に制限する必要がないときは「any」を使います。

\begin{verbatim}
$ sudo ufw insert 3 deny from 192.0.2.0/24 to any port 80 proto tcp
Rule inserted
$ sudo ufw status numbered
Status: active

     To                         Action      From
     --                         ------      ----
[ 1] 22/tcp                     ALLOW IN    Anywhere
[ 2] 53                         ALLOW IN    Anywhere
[ 3] 80/tcp                     DENY IN     192.0.2.0/24
[ 4] 80/tcp                     ALLOW IN    Anywhere
[ 5] 22/tcp (v6)                ALLOW IN    Anywhere (v6)
[ 6] 53 (v6)                    ALLOW IN    Anywhere (v6)
[ 7] 80/tcp (v6)                ALLOW IN    Anywhere (v6)

$ sudo ufw delete deny from 192.0.2.0/24 to any port 80 proto tcp
Rule deleted
$ sudo ufw insert 3 deny from 192.0.2.0/24 to any port 80 proto tcp
Rule inserted
$ sudo ufw delete 3
Deleting:
 deny from 192.0.2.0/24 to any port 80 proto tcp
Proceed with operation (y|n)? y
Rule deleted
\end{verbatim}

「from」で IPv4 アドレスを指定したので、ルールは自動的に IPv4 の方だけになって IPv6 には何も影響していません。

\subsection{on INTERFACE}

特定のネットワークインターフェイスを通るパケットだけ対象にしたい場合は「on INTERFACE」を使います。
例えばプロキシを必須にするなどで外部の 80,443 は直接繋がせない、とする場合は以下のようになります。
外部への通信なので「out」もつけています。

また、内部から外部への通信は「deny」で無視すると応答がなくなって、タイムアウト待ちなどでトラブルシューティングに時間がかかるので、すぐに拒否される「reject」を使うのがおすすめです。

インターフェイスは iptables と同様に「ppp+」で「ppp0」や「ppp1」などをまとめて指定できます。
そのため「enp0s3」だったり「ens3」だったりする stretch では「en+」でまとめて指定するとハードウェア構成に依存せずに設定を共有しやすくなると思います。

\begin{verbatim}
$ sudo ufw reject out on en+ to any port 80,443 proto tcp
Rule added
Rule added (v6)
$ sudo ufw status
Status: active

To                         Action      From
--                         ------      ----
22/tcp                     ALLOW       Anywhere
53                         ALLOW       Anywhere
80/tcp                     ALLOW       Anywhere
22/tcp (v6)                ALLOW       Anywhere (v6)
53 (v6)                    ALLOW       Anywhere (v6)
80/tcp (v6)                ALLOW       Anywhere (v6)

80,443/tcp                 REJECT OUT  Anywhere on en+
80,443/tcp (v6)            REJECT OUT  Anywhere (v6) on en+

$ curl --head www.debian.org
curl: (7) Failed to connect to www.debian.org port 80: Connection refused
$ sudo ufw delete reject out on en+ to any port 80,443 proto tcp
Rule deleted
Rule deleted (v6)
\end{verbatim}

deny だと応答がなくなったので Ctrl+C で止めました。

\begin{verbatim}
$ sudo ufw deny out on en+ to any port 80,443 proto tcp
Rule added
Rule added (v6)
$ curl --head www.debian.org
^C
$ sudo ufw delete deny out on en+ to any port 80,443 proto tcp
Rule deleted
Rule deleted (v6)
\end{verbatim}

\subsection{アプリケーションで指定}

「/etc/ufw/applications.d/」にアプリケーション名で設定するための情報が入っていて、「sudo ufw app list」で一覧できます。

\begin{verbatim}
$ sudo ufw app list
Available applications:
  AIM
  Bonjour
  CIFS
  DNS
  Deluge
  IMAP
  IMAPS
  IPP
  KTorrent
  Kerberos Admin
  Kerberos Full
  Kerberos KDC
  Kerberos Password
  LDAP
  LDAPS
  LPD
  MSN
  MSN SSL
  Mail submission
  NFS
  OpenSSH
  POP3
  POP3S
  PeopleNearby
  SMTP
  SSH
  Socks
  Telnet
  Transmission
  Transparent Proxy
  VNC
  WWW
  WWW Cache
  WWW Full
  WWW Secure
  XMPP
  Yahoo
  qBittorrent
  svnserve
$ ls /etc/ufw/applications.d/
openssh-server  ufw-chat             ufw-dnsserver   ufw-loginserver  ufw-printserver  ufw-webserver
ufw-bittorent   ufw-directoryserver  ufw-fileserver  ufw-mailserver   ufw-proxyserver
$ cat /etc/ufw/applications.d/openssh-server
[OpenSSH]
title=Secure shell server, an rshd replacement
description=OpenSSH is a free implementation of the Secure Shell protocol.
ports=22/tcp
$ cat /etc/ufw/applications.d/ufw-webserver
[WWW]
title=Web Server
description=Web server
ports=80/tcp

[WWW Secure]
title=Web Server (HTTPS)
description=Web Server (HTTPS)
ports=443/tcp

[WWW Full]
title=Web Server (HTTP,HTTPS)
description=Web Server (HTTP,HTTPS)
ports=80,443/tcp

[WWW Cache]
title=Web Server (8080)
description=Web Server (8080)
ports=8080/tcp
\end{verbatim}

アプリ名をポート番号とプロトコルの組み合わせの代わりに使うことができます。
スペースが入っているアプリ名の場合はシェルで引数が分割されないようにクォートする必要があるので、注意が必要です。

\begin{verbatim}
$ sudo ufw reject out OpenSSH
Rule added
Rule added (v6)
$ sudo ufw reject out to 10.0.0.0/8 app 'WWW Full'
Rule added
$ sudo ufw status
Status: active

To                         Action      From
--                         ------      ----
22/tcp                     ALLOW       Anywhere
53                         ALLOW       Anywhere
80/tcp                     ALLOW       Anywhere
22/tcp (v6)                ALLOW       Anywhere (v6)
53 (v6)                    ALLOW       Anywhere (v6)
80/tcp (v6)                ALLOW       Anywhere (v6)

OpenSSH                    REJECT OUT  Anywhere
10.0.0.0/8 WWW Full        REJECT OUT  Anywhere
OpenSSH (v6)               REJECT OUT  Anywhere (v6)
$ sudo ufw delete reject out to 10.0.0.0/8 app WWW Full
ERROR: Wrong number of arguments
$ sudo ufw delete reject out to 10.0.0.0/8 app 'WWW Full'
Rule deleted
$ sudo ufw delete reject out OpenSSH
Rule deleted
Rule deleted (v6)
\end{verbatim}

\subsection{limit}

拒否の仕方に deny と reject があるのと似た感じで、許可の方にも allow と limit の2種類があります。
limit は brute force attack の緩和に使えます。
30 秒間に 6 回までしか新規接続ができない、という設定になるようです。
正常な接続かどうかに関わらず、iptables の層で拒否してしまうので、例えば構成管理ツールなどで自分が短時間に頻繁に接続する可能性がある場合は避けた方が良いです。

\begin{verbatim}
$ sudo ufw limit 22/tcp
Rule updated
Rule updated (v6)
$ sudo ufw status
Status: active

To                         Action      From
--                         ------      ----
22/tcp                     LIMIT       Anywhere
53                         ALLOW       Anywhere
80/tcp                     ALLOW       Anywhere
22/tcp (v6)                LIMIT       Anywhere (v6)
53 (v6)                    ALLOW       Anywhere (v6)
80/tcp (v6)                ALLOW       Anywhere (v6)
\end{verbatim}

「/etc/ufw/user.rules」に設定されている内容は以下の通りです。

\begin{verbatim}
### tuple ### limit tcp 22 0.0.0.0/0 any 0.0.0.0/0 in
-A ufw-user-input -p tcp --dport 22 -m conntrack --ctstate NEW -m recent --set
-A ufw-user-input -p tcp --dport 22 -m conntrack --ctstate NEW -m recent --update \
   --seconds 30 --hitcount 6 -j ufw-user-limit (実際は1行)
-A ufw-user-input -p tcp --dport 22 -j ufw-user-limit-accept
\end{verbatim}

\subsection{コマンドラインで指定できない設定を入れる方法}

「/etc/ufw/before.rules」と「/etc/ufw/before6.rules」が ufw のルールの前に読み込まれる iptables-restore と ip6tables-restore のルールファイルになっているので、
IPsec の ESP の許可などの ufw コマンドで設定できないルールはこのファイルを直接編集すれば良いでしょう。
ufw で設定したルールより後に読み込まれる「after.rules」と「after6.rules」もあるので、用途によってはこちらを使っても良いでしょう。

\subsection{NAT 設定}

ちょっと複雑なネットワークだと、例えば POSTROUTING の MASQUERADE 設定はよく使うと思うのですが、「ufw」のコマンドラインでは nat テーブルの設定はできないので、「before.rules」などに追加することになります。

以下のような設定を「*filter」より上 (または COMMIT より下) に追加すると NAT の設定ができます。
反映させるには「ufw reload」を実行する必要があります。
また、「-F」の行が重要で、入れていないと「ufw reload」を実行するたびに nat テーブルの設定が増えていくことになります。

\begin{verbatim}
# NAT table rules
*nat
:POSTROUTING ACCEPT [0:0]
-F
# Allow traffic from OpenVPN client to enp0s3
-A POSTROUTING -s 192.168.10.0/24 -o en+ -j MASQUERADE
# Allow traffic from 192.168.1.0/24 (server's LAN subnet) to OpenVPN client
-A POSTROUTING -s 192.168.1.0/24 -o tun0 -j MASQUERADE
COMMIT
\end{verbatim}

iptables-restore で直接設定したものになるので、確認は「sudo iptables -t nat -nL」になります。

\begin{verbatim}
$ sudoedit /etc/ufw/before.rules
(-F がない場合)
$ sudo ufw reload
Firewall reloaded
$ sudo ufw reload
Firewall reloaded
$ sudo iptables -t nat -nL
Chain PREROUTING (policy ACCEPT)
target     prot opt source               destination

Chain INPUT (policy ACCEPT)
target     prot opt source               destination

Chain OUTPUT (policy ACCEPT)
target     prot opt source               destination

Chain POSTROUTING (policy ACCEPT)
target     prot opt source               destination
MASQUERADE  all  --  192.168.10.0/24      0.0.0.0/0
MASQUERADE  all  --  192.168.1.0/24       0.0.0.0/0
MASQUERADE  all  --  192.168.10.0/24      0.0.0.0/0
MASQUERADE  all  --  192.168.1.0/24       0.0.0.0/0
$ sudoedit /etc/ufw/before.rules
(-F を追加)
$ sudo ufw reload
Firewall reloaded
$ sudo iptables -t nat -nL
Chain PREROUTING (policy ACCEPT)
target     prot opt source               destination

Chain INPUT (policy ACCEPT)
target     prot opt source               destination

Chain OUTPUT (policy ACCEPT)
target     prot opt source               destination

Chain POSTROUTING (policy ACCEPT)
target     prot opt source               destination
MASQUERADE  all  --  192.168.10.0/24      0.0.0.0/0
MASQUERADE  all  --  192.168.1.0/24       0.0.0.0/0
\end{verbatim}

\subsection{routed の有効化}

初期状態では disabled になっていますが、「/etc/ufw/sysctl.conf」でも「/etc/sysctl.conf」でも「/etc/sysctl.d/*.conf」でも良いので、カーネルの設定を有効にすると「disabled」ではなくなります。

\begin{verbatim}
$ sudo ufw status verbose | grep Default
Default: deny (incoming), allow (outgoing), disabled (routed)
$ sudo tee /etc/sysctl.d/50-local.conf
net.ipv4.ip_forward=1
net.ipv6.conf.all.forwarding=1
$ sudo sysctl -p /etc/sysctl.d/50-local.conf
net.ipv4.ip_forward = 1
net.ipv6.conf.all.forwarding = 1
$ sudo ufw status verbose | grep Default
Default: deny (incoming), allow (outgoing), deny (routed)
\end{verbatim}

\subsection{FORWARD チェインの設定}

「ufw route」で iptables の FORWARD チェインの設定もできます。
「ufw route」は changelog によると ufw 0.34 から追加されているので、Debian だと jessie の 0.33-2 だと使えなくて stretch の 0.35-4 から使えるようです。

iptables の FORWARD チェインに対する設定になる以外は INPUT や OUTPUT チェインに対する操作と変わりません。

\begin{verbatim}
$ sudo ufw route allow from 192.168.10.0/24 to 10.0.0.0/8 port 80
Rule added
$ sudo ufw status
Status: active

To                         Action      From
--                         ------      ----
22/tcp                     LIMIT       Anywhere
53                         ALLOW       Anywhere
80/tcp                     ALLOW       Anywhere
22/tcp (v6)                LIMIT       Anywhere (v6)
53 (v6)                    ALLOW       Anywhere (v6)
80/tcp (v6)                ALLOW       Anywhere (v6)

10.0.0.0/8 80              ALLOW FWD   192.168.10.0/24
$ sudo ufw route delete allow from 192.168.10.0/24 to 10.0.0.0/8 port 80
Rule deleted
\end{verbatim}

\subsection{デフォルトポリシーの変更}

「 \verb~ufw [--dry-run] default allow|deny|reject [incoming|outgoing|routed]~ 」でチェーンのデフォルトのポリシーを変更できます。

stretch の ufw では、「default reject」は「 \verb~ufw*-reject-*~ 」チェーンで reject していて、FORWARD チェインなどの指定した組み込みのチェーンの policy が REJECT に変わるわけではないようです。

\begin{verbatim}
$ sudo ufw status verbose | grep Default
Default: deny (incoming), allow (outgoing), deny (routed)
$ grep FORWARD /etc/default/ufw
DEFAULT_FORWARD_POLICY="DROP"
$ sudo iptables -nL | grep FORWARD
Chain FORWARD (policy DROP)
$ sudo iptables -nL | grep -A2 'Chain.*reject-forward'
Chain ufw-reject-forward (1 references)
target     prot opt source               destination

$ sudo ufw default reject routed
Default routed policy changed to 'reject'
(be sure to update your rules accordingly)
$ grep FORWARD /etc/default/ufw
DEFAULT_FORWARD_POLICY="REJECT"
$ sudo iptables -nL | grep FORWARD
Chain FORWARD (policy DROP)
$ sudo ufw status verbose | grep Default
Default: deny (incoming), allow (outgoing), reject (routed)
$ sudo iptables -nL | grep -A2 'Chain.*reject-forward'
Chain ufw-reject-forward (1 references)
target     prot opt source               destination
REJECT     all  --  0.0.0.0/0            0.0.0.0/0            reject-with icmp-port-unreachable
$ sudo ufw default deny routed
Default routed policy changed to 'deny'
(be sure to update your rules accordingly)
\end{verbatim}

\subsection{コメント}

今の ufw のルールにはコメントがつけられるようになっています。
コメントは ufw のルールに紐づいていて、 iptables の方には特に何も影響はないようです。

\begin{verbatim}
$ sudo ufw allow 60000:61000/udp comment "mobile shell"
Rule added
Rule added (v6)
$ sudo ufw status | grep 60000
60000:61000/udp            ALLOW       Anywhere                   # mobile shell
60000:61000/udp (v6)       ALLOW       Anywhere (v6)              # mobile shell
$ sudo grep 60000 /etc/ufw/user*.rules
/etc/ufw/user6.rules:### tuple ### allow udp 60000:61000 ::/0 any ::/0 in comment=6d6f62696c65207368656c6c
/etc/ufw/user6.rules:-A ufw6-user-input -p udp -m multiport --dports 60000:61000 -j ACCEPT
/etc/ufw/user.rules:### tuple ### allow udp 60000:61000 0.0.0.0/0 any 0.0.0.0/0 in comment=6d6f62696c65207368656c6c
/etc/ufw/user.rules:-A ufw-user-input -p udp -m multiport --dports 60000:61000 -j ACCEPT
\end{verbatim}

\subsection{ufw show REPORT}

ufw の man の REPORTS セクションで紹介されているように、「ufw show REPORT」で設定状態などを表示できます。
「sudo ufw show listening」は実際に待ち受けしているポートを ufw の設定状況と一緒に一覧できるので、便利そうです。

\begin{itemize}
\item sudo ufw show raw
\item sudo ufw show builtins
\item sudo ufw show before-rules
\item sudo ufw show user-rules
\item sudo ufw show after-rules
\item sudo ufw show logging-rules
\item sudo ufw show listening
\item sudo ufw show added
\end{itemize}

\begin{verbatim}
$ sudo ufw show listening
tcp:
  111 * (rpcbind)
  22 * (sshd)
   [ 1] allow 22/tcp

tcp6:
  111 * (rpcbind)
  22 * (sshd)
   [ 2] allow 22/tcp

udp:
  1015 * (rpcbind)
  111 * (rpcbind)
  68 * (dhclient)
udp6:
  1015 * (rpcbind)
  111 * (rpcbind)
\end{verbatim}

\subsection{ログレベル}

デフォルトだと拒否された時のログが rate limit 付きで /var/log/ufw.log などに残りますが、

\begin{verbatim}
ufw [--dry-run] logging on|off|LEVEL
\end{verbatim}

で、どのくらいログを残すのか設定できます。

設定できるのは on, off, low, medium, high, full のようです。
詳細は man を参照してください。


\dancersection{今後の予定}{Debian JP}

\subsection{関西Debian勉強会}
次回、第133回関西Debian勉強会は3月25(日)に開催予定です。

\subsection{東京エリアDebian勉強会}
次回、第161回東京エリアDebian勉強会は3月24日(土)に新宿で開催予定です。
内容は機械学習環境構築の予定です。

%
% 冊子にするために、4の倍数にする必要がある。
% そのための調整
\dancersection{メモ}{}
\mbox{}\newpage
% \mbox{}\newpage
% \mbox{}\newpage

\pagebreak

\begin{center}
本資料のライセンスについて
\end{center}

本資料はフリー・ソフトウェアです。あなたは、Free Software
Foundation が公表したGNU GENERAL PUBLIC LICENSEの "バージョン２"もしくはそれ以降
が定める条項に従って本プログラムを再頒布または変更することができ
ます。

本プログラムは有用とは思いますが、頒布にあたっては、市場性及び特
定目的適合性についての暗黙の保証を含めて、いかなる保証も行ないま
せん。詳細についてはGNU GENERAL PUBLIC LICENSE をお読みください。

\begin{multicols}{2}
 \begin{fontsize}{6}{6}
 \begin{verbatim}
            GNU GENERAL PUBLIC LICENSE
               Version 2, June 1991

 Copyright (C) 1989, 1991 Free Software Foundation, Inc.
    51 Franklin St, Fifth Floor, Boston, MA  02110-1301  USA
 Everyone is permitted to copy and distribute verbatim copies
 of this license document, but changing it is not allowed.

                Preamble

  The licenses for most software are designed to take away your
freedom to share and change it.  By contrast, the GNU General Public
License is intended to guarantee your freedom to share and change free
software--to make sure the software is free for all its users.  This
General Public License applies to most of the Free Software
Foundation's software and to any other program whose authors commit to
using it.  (Some other Free Software Foundation software is covered by
the GNU Library General Public License instead.)  You can apply it to
your programs, too.

  When we speak of free software, we are referring to freedom, not
price.  Our General Public Licenses are designed to make sure that you
have the freedom to distribute copies of free software (and charge for
this service if you wish), that you receive source code or can get it
if you want it, that you can change the software or use pieces of it
in new free programs; and that you know you can do these things.

  To protect your rights, we need to make restrictions that forbid
anyone to deny you these rights or to ask you to surrender the rights.
These restrictions translate to certain responsibilities for you if you
distribute copies of the software, or if you modify it.

  For example, if you distribute copies of such a program, whether
gratis or for a fee, you must give the recipients all the rights that
you have.  You must make sure that they, too, receive or can get the
source code.  And you must show them these terms so they know their
rights.

  We protect your rights with two steps: (1) copyright the software, and
(2) offer you this license which gives you legal permission to copy,
distribute and/or modify the software.

  Also, for each author's protection and ours, we want to make certain
that everyone understands that there is no warranty for this free
software.  If the software is modified by someone else and passed on, we
want its recipients to know that what they have is not the original, so
that any problems introduced by others will not reflect on the original
authors' reputations.

  Finally, any free program is threatened constantly by software
patents.  We wish to avoid the danger that redistributors of a free
program will individually obtain patent licenses, in effect making the
program proprietary.  To prevent this, we have made it clear that any
patent must be licensed for everyone's free use or not licensed at all.

  The precise terms and conditions for copying, distribution and
modification follow.

            GNU GENERAL PUBLIC LICENSE
   TERMS AND CONDITIONS FOR COPYING, DISTRIBUTION AND MODIFICATION

  0. This License applies to any program or other work which contains
a notice placed by the copyright holder saying it may be distributed
under the terms of this General Public License.  The "Program", below,
refers to any such program or work, and a "work based on the Program"
means either the Program or any derivative work under copyright law:
that is to say, a work containing the Program or a portion of it,
either verbatim or with modifications and/or translated into another
language.  (Hereinafter, translation is included without limitation in
the term "modification".)  Each licensee is addressed as "you".

Activities other than copying, distribution and modification are not
covered by this License; they are outside its scope.  The act of
running the Program is not restricted, and the output from the Program
is covered only if its contents constitute a work based on the
Program (independent of having been made by running the Program).
Whether that is true depends on what the Program does.

  1. You may copy and distribute verbatim copies of the Program's
source code as you receive it, in any medium, provided that you
conspicuously and appropriately publish on each copy an appropriate
copyright notice and disclaimer of warranty; keep intact all the
notices that refer to this License and to the absence of any warranty;
and give any other recipients of the Program a copy of this License
along with the Program.

You may charge a fee for the physical act of transferring a copy, and
you may at your option offer warranty protection in exchange for a fee.

  2. You may modify your copy or copies of the Program or any portion
of it, thus forming a work based on the Program, and copy and
distribute such modifications or work under the terms of Section 1
above, provided that you also meet all of these conditions:

    a) You must cause the modified files to carry prominent notices
    stating that you changed the files and the date of any change.

    b) You must cause any work that you distribute or publish, that in
    whole or in part contains or is derived from the Program or any
    part thereof, to be licensed as a whole at no charge to all third
    parties under the terms of this License.

    c) If the modified program normally reads commands interactively
    when run, you must cause it, when started running for such
    interactive use in the most ordinary way, to print or display an
    announcement including an appropriate copyright notice and a
    notice that there is no warranty (or else, saying that you provide
    a warranty) and that users may redistribute the program under
    these conditions, and telling the user how to view a copy of this
    License.  (Exception: if the Program itself is interactive but
    does not normally print such an announcement, your work based on
    the Program is not required to print an announcement.)

These requirements apply to the modified work as a whole.  If
identifiable sections of that work are not derived from the Program,
and can be reasonably considered independent and separate works in
themselves, then this License, and its terms, do not apply to those
sections when you distribute them as separate works.  But when you
distribute the same sections as part of a whole which is a work based
on the Program, the distribution of the whole must be on the terms of
this License, whose permissions for other licensees extend to the
entire whole, and thus to each and every part regardless of who wrote it.

Thus, it is not the intent of this section to claim rights or contest
your rights to work written entirely by you; rather, the intent is to
exercise the right to control the distribution of derivative or
collective works based on the Program.

In addition, mere aggregation of another work not based on the Program
with the Program (or with a work based on the Program) on a volume of
a storage or distribution medium does not bring the other work under
the scope of this License.

  3. You may copy and distribute the Program (or a work based on it,
under Section 2) in object code or executable form under the terms of
Sections 1 and 2 above provided that you also do one of the following:

    a) Accompany it with the complete corresponding machine-readable
    source code, which must be distributed under the terms of Sections
    1 and 2 above on a medium customarily used for software interchange; or,

    b) Accompany it with a written offer, valid for at least three
    years, to give any third party, for a charge no more than your
    cost of physically performing source distribution, a complete
    machine-readable copy of the corresponding source code, to be
    distributed under the terms of Sections 1 and 2 above on a medium
    customarily used for software interchange; or,

    c) Accompany it with the information you received as to the offer
    to distribute corresponding source code.  (This alternative is
    allowed only for noncommercial distribution and only if you
    received the program in object code or executable form with such
    an offer, in accord with Subsection b above.)

The source code for a work means the preferred form of the work for
making modifications to it.  For an executable work, complete source
code means all the source code for all modules it contains, plus any
associated interface definition files, plus the scripts used to
control compilation and installation of the executable.  However, as a
special exception, the source code distributed need not include
anything that is normally distributed (in either source or binary
form) with the major components (compiler, kernel, and so on) of the
operating system on which the executable runs, unless that component
itself accompanies the executable.

If distribution of executable or object code is made by offering
access to copy from a designated place, then offering equivalent
access to copy the source code from the same place counts as
distribution of the source code, even though third parties are not
compelled to copy the source along with the object code.

  4. You may not copy, modify, sublicense, or distribute the Program
except as expressly provided under this License.  Any attempt
otherwise to copy, modify, sublicense or distribute the Program is
void, and will automatically terminate your rights under this License.
However, parties who have received copies, or rights, from you under
this License will not have their licenses terminated so long as such
parties remain in full compliance.

  5. You are not required to accept this License, since you have not
signed it.  However, nothing else grants you permission to modify or
distribute the Program or its derivative works.  These actions are
prohibited by law if you do not accept this License.  Therefore, by
modifying or distributing the Program (or any work based on the
Program), you indicate your acceptance of this License to do so, and
all its terms and conditions for copying, distributing or modifying
the Program or works based on it.

  6. Each time you redistribute the Program (or any work based on the
Program), the recipient automatically receives a license from the
original licensor to copy, distribute or modify the Program subject to
these terms and conditions.  You may not impose any further
restrictions on the recipients' exercise of the rights granted herein.
You are not responsible for enforcing compliance by third parties to
this License.

  7. If, as a consequence of a court judgment or allegation of patent
infringement or for any other reason (not limited to patent issues),
conditions are imposed on you (whether by court order, agreement or
otherwise) that contradict the conditions of this License, they do not
excuse you from the conditions of this License.  If you cannot
distribute so as to satisfy simultaneously your obligations under this
License and any other pertinent obligations, then as a consequence you
may not distribute the Program at all.  For example, if a patent
license would not permit royalty-free redistribution of the Program by
all those who receive copies directly or indirectly through you, then
the only way you could satisfy both it and this License would be to
refrain entirely from distribution of the Program.

If any portion of this section is held invalid or unenforceable under
any particular circumstance, the balance of the section is intended to
apply and the section as a whole is intended to apply in other
circumstances.

It is not the purpose of this section to induce you to infringe any
patents or other property right claims or to contest validity of any
such claims; this section has the sole purpose of protecting the
integrity of the free software distribution system, which is
implemented by public license practices.  Many people have made
generous contributions to the wide range of software distributed
through that system in reliance on consistent application of that
system; it is up to the author/donor to decide if he or she is willing
to distribute software through any other system and a licensee cannot
impose that choice.

This section is intended to make thoroughly clear what is believed to
be a consequence of the rest of this License.

  8. If the distribution and/or use of the Program is restricted in
certain countries either by patents or by copyrighted interfaces, the
original copyright holder who places the Program under this License
may add an explicit geographical distribution limitation excluding
those countries, so that distribution is permitted only in or among
countries not thus excluded.  In such case, this License incorporates
the limitation as if written in the body of this License.

  9. The Free Software Foundation may publish revised and/or new versions
of the General Public License from time to time.  Such new versions will
be similar in spirit to the present version, but may differ in detail to
address new problems or concerns.

Each version is given a distinguishing version number.  If the Program
specifies a version number of this License which applies to it and "any
later version", you have the option of following the terms and conditions
either of that version or of any later version published by the Free
Software Foundation.  If the Program does not specify a version number of
this License, you may choose any version ever published by the Free Software
Foundation.

  10. If you wish to incorporate parts of the Program into other free
programs whose distribution conditions are different, write to the author
to ask for permission.  For software which is copyrighted by the Free
Software Foundation, write to the Free Software Foundation; we sometimes
make exceptions for this.  Our decision will be guided by the two goals
of preserving the free status of all derivatives of our free software and
of promoting the sharing and reuse of software generally.

                NO WARRANTY

  11. BECAUSE THE PROGRAM IS LICENSED FREE OF CHARGE, THERE IS NO WARRANTY
FOR THE PROGRAM, TO THE EXTENT PERMITTED BY APPLICABLE LAW.  EXCEPT WHEN
OTHERWISE STATED IN WRITING THE COPYRIGHT HOLDERS AND/OR OTHER PARTIES
PROVIDE THE PROGRAM "AS IS" WITHOUT WARRANTY OF ANY KIND, EITHER EXPRESSED
OR IMPLIED, INCLUDING, BUT NOT LIMITED TO, THE IMPLIED WARRANTIES OF
MERCHANTABILITY AND FITNESS FOR A PARTICULAR PURPOSE.  THE ENTIRE RISK AS
TO THE QUALITY AND PERFORMANCE OF THE PROGRAM IS WITH YOU.  SHOULD THE
PROGRAM PROVE DEFECTIVE, YOU ASSUME THE COST OF ALL NECESSARY SERVICING,
REPAIR OR CORRECTION.

  12. IN NO EVENT UNLESS REQUIRED BY APPLICABLE LAW OR AGREED TO IN WRITING
WILL ANY COPYRIGHT HOLDER, OR ANY OTHER PARTY WHO MAY MODIFY AND/OR
REDISTRIBUTE THE PROGRAM AS PERMITTED ABOVE, BE LIABLE TO YOU FOR DAMAGES,
INCLUDING ANY GENERAL, SPECIAL, INCIDENTAL OR CONSEQUENTIAL DAMAGES ARISING
OUT OF THE USE OR INABILITY TO USE THE PROGRAM (INCLUDING BUT NOT LIMITED
TO LOSS OF DATA OR DATA BEING RENDERED INACCURATE OR LOSSES SUSTAINED BY
YOU OR THIRD PARTIES OR A FAILURE OF THE PROGRAM TO OPERATE WITH ANY OTHER
PROGRAMS), EVEN IF SUCH HOLDER OR OTHER PARTY HAS BEEN ADVISED OF THE
POSSIBILITY OF SUCH DAMAGES.

             END OF TERMS AND CONDITIONS

        How to Apply These Terms to Your New Programs

  If you develop a new program, and you want it to be of the greatest
possible use to the public, the best way to achieve this is to make it
free software which everyone can redistribute and change under these terms.

  To do so, attach the following notices to the program.  It is safest
to attach them to the start of each source file to most effectively
convey the exclusion of warranty; and each file should have at least
the "copyright" line and a pointer to where the full notice is found.

    <one line to give the program's name and a brief idea of what it does.>
    Copyright (C) <year>  <name of author>

    This program is free software; you can redistribute it and/or modify
    it under the terms of the GNU General Public License as published by
    the Free Software Foundation; either version 2 of the License, or
    (at your option) any later version.

    This program is distributed in the hope that it will be useful,
    but WITHOUT ANY WARRANTY; without even the implied warranty of
    MERCHANTABILITY or FITNESS FOR A PARTICULAR PURPOSE.  See the
    GNU General Public License for more details.

    You should have received a copy of the GNU General Public License
    along with this program; if not, write to the Free Software
    Foundation, Inc., 51 Franklin St, Fifth Floor, Boston, MA  02110-1301 USA


Also add information on how to contact you by electronic and paper mail.

If the program is interactive, make it output a short notice like this
when it starts in an interactive mode:

    Gnomovision version 69, Copyright (C) year  name of author
    Gnomovision comes with ABSOLUTELY NO WARRANTY; for details type `show w'.
    This is free software, and you are welcome to redistribute it
    under certain conditions; type `show c' for details.

The hypothetical commands `show w' and `show c' should show the appropriate
parts of the General Public License.  Of course, the commands you use may
be called something other than `show w' and `show c'; they could even be
mouse-clicks or menu items--whatever suits your program.

You should also get your employer (if you work as a programmer) or your
school, if any, to sign a "copyright disclaimer" for the program, if
necessary.  Here is a sample; alter the names:

  Yoyodyne, Inc., hereby disclaims all copyright interest in the program
  `Gnomovision' (which makes passes at compilers) written by James Hacker.

  <signature of Ty Coon>, 1 April 1989
  Ty Coon, President of Vice

This General Public License does not permit incorporating your program into
proprietary programs.  If your program is a subroutine library, you may
consider it more useful to permit linking proprietary applications with the
library.  If this is what you want to do, use the GNU Library General
Public License instead of this License.
 \end{verbatim}
 \end{fontsize}
\end{multicols}

\begin{center}
ソースコードについて
\end{center}

このプログラムは tex で記述されたものです。ソースコードは
\begin{center}
  \url{git://anonscm.debian.org/tokyodebian/monthly-report.git}
\end{center}
から取得できます。

\begin{center}
Debian オープンユーズロゴ ライセンス
\end{center}

\begin{multicols}{2}
 \begin{fontsize}{6}{6}
 \begin{verbatim}

Copyright (c) 1999 Software in the Public Interest
Permission is hereby granted, free of charge, to any person
obtaining a copy of this software and associated documentation
files (the "Software"), to deal in the Software without restriction,
including without limitation the rights to use, copy, modify, merge,
publish, distribute, sublicense, and/or sell copies of the Software,
and to permit persons to whom the Software is furnished to do so,
subject to the following conditions:

The above copyright notice and this permission notice shall be
included in all copies or substantial portions of the Software.

THE SOFTWARE IS PROVIDED "AS IS", WITHOUT WARRANTY OF ANY
KIND, EXPRESS OR IMPLIED, INCLUDING BUT NOT LIMITED TO THE
WARRANTIES OF MERCHANTABILITY, FITNESS FOR A PARTICULAR PURPOSE AND
NONINFRINGEMENT. IN NO EVENT SHALL THE AUTHORS OR COPYRIGHT HOLDERS
BE LIABLE FOR ANY CLAIM, DAMAGES OR OTHER LIABILITY, WHETHER IN
AN ACTION OF CONTRACT, TORT OR OTHERWISE, ARISING FROM, OUT OF OR
IN CONNECTION WITH THE SOFTWARE OR THE USE OR OTHER DEALINGS IN
THE SOFTWARE.
 \end{verbatim}
 \end{fontsize}
\end{multicols}

\printindex
%\cleartooddpage

 \begin{minipage}[b]{0.2\hsize}
  \rotatebox{90}{\fontsize{80}{80} {\gt 関西 Debian 勉強会} }
 \end{minipage}
 \begin{minipage}[b]{0.8\hsize}

 \vspace*{15cm}
 \rule{\hsize}{1mm}
 \vspace{2mm}
 \includegraphics[width=2cm]{image200502/openlogo-nd.eps}
 \noindent \Large \bfseries{Debian 勉強会資料}\\ \\
 \noindent \normalfont \debmtgyear{}年\debmtgmonth{}月\debmtgdate{}日 \hspace{5mm}  初版第1刷発行\\
 \noindent \normalfont 関西 Debian 勉強会 （編集・印刷・発行）\\
 \rule{\hsize}{1mm}
 \end{minipage}

\end{document}
%%% Local Variables:
%%% skk-kutouten-type: jp
%%% End:
